%%%PAGE 186 rot=0 legibility=fair summary=实验笔记：验证力的平行四边形定则（钩码三绳）、探究动能定理（斜面装置与小车装置）、弹簧弹性势能实验（光电门、砝码测k）、验证机械能守恒（自由落体）
% 说明：本页是粘贴印刷纸被掀开后的笔记页全貌（p185、p187 为同一笔记页、被不同印刷纸盖住中段的另外两张照片，其上沿/下沿的相同笔迹只在本页转录）。
% 页首右上角另有一段淡铅笔化学字迹（n CO2 … 燃烧得CO和CO2混合气体 ρ=1.429 g·L^-1；CO为0.4×3/4×22.4=6.72 L），
% 在 p189 的同一位置也原样出现，应是垫在笔记本下面的另一张纸/页的透印，不是本页内容，故未作为本页的块转录（保留于此备查）。
%%%BLOCK id=186-01 ch=02 sec=验证力的平行四边形定则 kind=note alt=none cont=none y=0.03-0.21 exp=1
%FIG p186_a 0.04 0.027 0.22 0.165 soft
\notefig[0.25]{figures/p186_a.png}

要记录结点 $O$ 的位置、每条绳上所挂钩码的个数，$OA$、$OB$、$OC$\unclear{…}（行间小字：大小方向）
% 此行右侧为淡铅笔字迹、对焦模糊：
（右侧淡字：\unclear{分力的大小、方向作用……}；\unclear{$F_B$ 沿绳的方向}）

不用让 $OA\perp OB$，不用让 $O$ 点不变，可以换成弹性绳

（只需稳定即可）\quad（分力和合力可在同一个图示中表示出来）\quad（不影响受力的方向和\unclear{钩码}个数）
%%%END
%%%BLOCK id=186-02 ch=06 sec=探究动能定理·斜面装置 kind=note alt=none cont=none y=0.19-0.40 exp=1
%FIG p186_b 0.04 0.19 0.30 0.29 soft
\notefig[0.25]{figures/p186_b.png}

测 $v$ 则不用平衡了 $f$

$W_{\text{合}}=mgx=\dfrac12(M+m)v^2$，\quad $v^2=\dfrac{2}{M+m}W$ 是一条过原点的倾斜直线

$x-v^2$ \\
$W-v^2$ \\
$v^2-W$

还可由 $k$ 求出滑块质量 $M$
%%%END
%%%BLOCK id=186-03 ch=06 sec=验证机械能守恒定律 kind=note alt=none cont=none y=0.33-0.38 exp=1
单体机守不用测 $m$\quad 理由：物体质量可约掉
%%%END
%%%BLOCK id=186-04 ch=06 sec=探究弹簧弹性势能·光电门 kind=note alt=none cont=none y=0.40-0.58 exp=1
$\dfrac12k(\Delta x)^2=\dfrac12mv^2$，\quad $\Delta x\propto v$

%FIG p186_c 0.46 0.413 0.74 0.463
\notefig[0.28]{figures/p186_c.png}

\begin{keybox}
对同一根弹簧：弹性势能与弹簧的压缩量平方成正比。
\end{keybox}

弹簧弹性势能转化为滑块动能

\begin{center}
\begin{tabular}{@{}lccc@{}}
砝码质量$/\unit{g}$ & 50 & 100 & 150 \\
 & 8.62 & 7.63 & 6.66 \\
\end{tabular}
\end{center}
% CHECK: 第二行数据（8.62、7.63、6.66）原稿无行名，疑为弹簧长度/cm
（8.62 与 7.63 之间画有双向箭头，其下标 $\Delta x=0.99$；其后 $\Delta x=0.97$）\quad 取平均 $\Delta x=0.98$

$k=\dfrac{\Delta f}{\Delta x}=\dfrac{50\times10^{-3}\times9.8}{0.98\times10^{-2}}=50\unit{N/m}$% CHECK: 分子可能是 ΔF
%%%END
%%%BLOCK id=186-05 ch=06 sec=探究动能定理·小车装置 kind=note alt=02 cont=none y=0.55-0.74 exp=1
%FIG p186_d 0.03 0.555 0.46 0.66
\notefig[0.43]{figures/p186_d.png}

探究动能定理\quad $T=\dfrac{Mmg}{M+m}=\dfrac{mg}{1+\frac{m}{M}}\approx mg$

（把\unclear{砂和砂}桶总重力当作小车受到的合外力.）\quad 要平衡摩擦，$m\ll M$

%FIG p186_e 0.12 0.655 0.53 0.732
\notefig[0.42]{figures/p186_e.png}

（图下标注：↑ 间隔3个点）

刻度尺（箭头指向 $x$）；天平（两箭头分别指向 $mg$ 与 $M$）

$mgx=\dfrac12M\left[\left(\dfrac{x_2}{4T}\right)^2-\left(\dfrac{x_1}{4T}\right)^2\right]$

间隔 $n$ 个点\quad $v=\dfrac{x_1+x_2}{2(n+1)T}$
%%%END
%%%BLOCK id=186-06 ch=06 sec=验证机械能守恒定律 kind=note alt=none cont=none y=0.74-0.97 exp=1
% 页脚右侧有一小段涂鸦弧线，忽略
自由落体验机守

%FIG p186_f 0.02 0.77 0.22 0.85
\notefig[0.25]{figures/p186_f.png}

有两处不当：① 用交流电不用直流电；② 重物离打点计时器太远。

利用第 $n$ 点与起始点计算：$v^2=2gh$，$\dfrac12v^2=gh$，图线斜率表示重力加速度 $g$
%%%END
%%%PAGEEND 186
